% V18 integration of Dot batch 097; immutable source SHA256 93a90e26edcb51ba4626bd690439f687907dcd74e932da78000a7eae3745a938
\begingroup
% These definitions have only group-local scope; the shared preamble is untouched.
\def\OO{\mathcal O}\def\FF{\mathcal F}\def\GG{\mathcal G}
\def\PP{\mathbf P}\def\calP{\mathcal P}\def\dd{\delta}
\def\Div{\operatorname{Div}}\def\depth{\operatorname{depth}}
\def\pr{\operatorname{pr}}\def\coker{\operatorname{coker}}
\egasection{§7\quad 应用于构造指定类型的超平面截面或多截面}

\egap{7.1}
设 $S=\Spec k$，且 $k$ 是无限域。那么 $P^\vee$ 的任一非空开子集都含有 $k$-有理点。采用4.1的记号，若由可构造性质 $\calP$ 定出的 $E$ 含泛点 $\eta$，则 $E$ 含一个 $k$-有理点。因此，$X$ 存在一个在 $k$ 上定义且满足 $\calP$ 的超平面截面。

现在让 $S$ 仍为任意基底。给定 $s\in S$ 及纤维 $P^\vee_s$ 的一个 $k(s)$-有理点 $\xi_0$，可以在 $s$ 的一个开邻域 $U$ 上找到 $P^\vee\to S$ 的截面 $\xi$，使它经过 $\xi_0$。仍如4.1以可构造性质 $\calP$ 定义 $E$。若已知 $E$ 是开的（例如由第5节得到这一点），且 $\xi_0\in E$，则充分缩小 $U$ 后，$\xi$ 的像落在 $E$ 中。于是得到 $X$ 在 $U$ 上的一个超平面截面 $Y_\xi$，使每个 $t\in U$ 上的纤维 $Y_{\xi(t)}$ 都满足 $\calP$。\footnote{法26将初始点手改成 $\xi_0$、将开放性的依据改指第5节，并补记 $Y_\xi$；法27首行手补 $Y_{\xi(t)}$。区别参数丛的截面 $\xi$ 和 $X$ 的超平面截面 $Y_\xi$。}

若事先没有指定 $\xi_0$，而 $k(s)$ 是无限域，可把前两点结合起来，在 $s$ 的某个开邻域上构造具有上述性质的超平面截面。\footnote{这里仍须在 $P^\vee_s$ 中有一个非空的好参数开集，例如纤维内相应好参数集可构造且含其泛点。单凭 $k(s)$ 无限，不能推出有好截面。}
最后，利用第5节可得到如下判据：在 $X$（相应地 $\FF$）相对 $S$ 平坦并满足该节判据的其余条件时，$Y_\xi$（相应地 $\GG_\xi$）也相对 $S$ 平坦。因而可用 $Y_\xi$ 代替 $X$，重复上述构造。比如，在一定条件下，可以逐步构造 $X$ 在给定点 $s$ 的某个开邻域 $U$ 上的一条“多截面” $S'$，使 $S'\to U$ 有限、平坦且各纤维满足 $\calP$。\footnote{“逐步构造”取法文清楚可辨的 \emph{de proche en proche}，不将英译问号处译作另有极限的“逐次逼近”。法稿说“例如，在一定条件下”，未在这里给出无条件多截面定理。}

若 $k(s)$ 是有限域，可能须先作满射的 étale 换底 $S'\to U$，再作上述构造。确实，与7.1开头相同，考虑 $P^\vee$ 中的非空开集；若 $k$ 有限，该开集未必有 $k$-有理点，但必有闭点，给出取值于 $k$ 的有限扩域 $k'$ 的点。这个扩域必可分。令 $k=k(s)$，在 $s$ 的某个开邻域 $U$ 上作与剩余域扩张 $k'/k(s)$ 对应的适当有限 étale 覆盖 $S'\to U$，使
\[
 S'_s\simeq\Spec k'.
\]
于是经换底 $S'\to S$，可在 $s$ 上方唯一的点 $s'\in S'$ 处归结到有一个有理好参数点的情形。\footnote{法27、英19均写“6.1开头”，实际说的是本条7.1起首的非空开集与有理点。该段把参数空间的非空开集和基底邻域都写成 $U$，译文分开。有限扩域并不会变成无限域；有利之处是选定的好闭点在新剩余域上成为有理点。这里的 $S'$ 表示换底覆盖，不是上一段的目标多截面。}
下面给出用这些方法得到的一些命题的例子。

\par\noindent\textit{译校补证另见核查报告。}\cmt{D18-097-R}{无限域选点、局部坐标提升及有限平展覆盖的补证完整保存在本版核查报告“D18-097”条；其中译校补证为非原文。}
\subsection*{法文划销旧稿的校记（非定稿）}
\paragraph{法26上部：关于第4节的改写建议.}
作者说，旧4.6(b)、(c)的结论已包含于旧4.7(b)、(c)。起初因后者的证明似乎需要更细致的结果（用15.2.2的形式），才认为从证明角度可以保留重复。若只想证明相应集合 $U$（分别为 $E$）开，而不证明 $\GG|_U$ 相对 $P^\vee$ 平坦，可以改走第12段使用的赋值判据。作者建议删去4.6(b)、(c)，将4.7(b)、(c)仅表述成开性（同时删4.7(a)），再添命题或推论单列平坦性并用15.2.2。行间可见手记“4.8”。全段被划销，编号均按废稿记，不据此改写现行第4、5节。正式 EGA~IV 第三分册（15.2.2）确是带附加条件的平坦性定理，见\href{https://www.numdam.org/item/PMIHES_1966__28__5_0/}{印刷页226}；不是“开便平坦”。\footnote{废稿局部断字，以上以能辨认的三条编辑建议及当时编号为准。正式（15.2.2）要求支撑、相应泛点处的普遍开性、模纤维几何约化和底局部环约化等条件；未宣称废稿对这些条件的应用已经完证。}

\paragraph{法27下部：与旧4.6、4.7相关的另一则剪裁稿.}
“下面给出……例子”一句未删。其后被划销的片段准备讨论如何放宽涉及超平面和 $\Supp\FF_s$ 的几何不可约分支的假设，拟给 $P^\vee$ 取一个开子集 $W$ 后改写条件。右端和末句不全，本批只登记可辨方向，不据片段补造完整假设。法28起笔是另一个未删的关于4.2与4.3的补记，接在下文。

\endgroup
